\documentclass[11pt,a4paper]{article}

\usepackage[margin=1in]{geometry}
\usepackage[T1]{fontenc}
\usepackage{graphicx}
\usepackage{booktabs}
\usepackage{amsmath}
\usepackage{hyperref}
\usepackage{enumitem}
\usepackage{float}
\usepackage{tabularx}

\graphicspath{{images/}}
\urlstyle{same}
\setlist[itemize]{leftmargin=1.4em,itemsep=0.25em,topsep=0.3em}

\title{\textbf{AI for Mental Health: Credible Signal Detection System}}
\author{Ujjwal Singh \\ \small ujjosing@gmail.com \\ \small \href{https://github.com/fallofpheonix}{github.com/fallofpheonix}}
\date{}

\begin{document}

\maketitle

\section{Abstract}

Social media can provide early signs of mental health crises, but the signals are often weak, noisy, and easy to distort. Most existing systems focus on classification accuracy, but in practice the more important question is whether an alert is reliable enough to act on. False alerts can waste limited human attention, while weak but real signals in low-data regions can be missed.

This proposal introduces a credible signal detection system that measures both crisis intensity and signal reliability. The pipeline combines data ingestion, signal extraction, credibility modeling, threat filtering, and a confidence-gated decision engine. High-risk cases are sent for human review instead of automatic action.

The expected outcome is a simpler and more trustworthy monitoring pipeline that can produce interpretable alerts under sparse, noisy, and adversarial conditions.

\section{Problem Statement}

Current crisis-monitoring systems usually work like direct classification pipelines. They ingest public posts, score them, group them by region or time, and escalate if the score crosses a threshold. This design breaks in real settings for three reasons:

\begin{itemize}
\item \textbf{Sparse data:} low-activity regions produce unstable estimates and small-number artifacts.
\item \textbf{Adversarial or media-driven noise:} coordinated amplification or news bursts can inflate apparent severity without reflecting genuine local distress.
\item \textbf{No explicit uncertainty:} most systems do not represent whether an alert is statistically credible enough to trigger review.
\end{itemize}

Figure \ref{fig:current-system} shows a typical direct pipeline where signal strength flows into escalation without a separate reliability check.

\begin{figure}[htbp]
\centering
\includegraphics[width=0.88\textwidth]{current_system.png}
\caption{Typical crisis-monitoring pipeline without explicit confidence calibration.}
\label{fig:current-system}
\end{figure}

\vspace{0.5cm}

The core problem is not only whether crisis-related content exists, but whether the observed evidence supports a \emph{credible} signal:

\[
P(\text{credible signal} \mid \text{sparse, noisy, adversarial data})
\]

\section{Proposed Solution}

The proposed system treats crisis monitoring as a credibility-aware inference problem. Instead of escalating directly from base signal estimates, it computes a final credible-signal score that includes both statistical support and adversarial suppression. This separates \emph{severity} from \emph{trustworthiness}, which is essential for real use.

Figure \ref{fig:proposed-architecture} presents the proposed governance-ready architecture. Relative to the baseline pipeline, it adds two explicit components: credibility modeling and threat filtering.

\begin{figure}[htbp]
\centering
\includegraphics[width=0.9\textwidth]{proposed_architecture.png}
\caption{Proposed governance-ready architecture for credible signal detection.}
\label{fig:proposed-architecture}
\end{figure}

\vspace{0.5cm}

As shown in Figure \ref{fig:proposed-architecture}, the design routes raw posts through preprocessing, signal extraction, credibility modeling, adversarial filtering, and confidence-gated escalation. Any high-priority output still goes through human review and audit logging.

\section{System Architecture}

\subsection{Overview}

The system is organized into five principal computational stages and one governance stage:

\begin{itemize}
\item \textbf{Data ingestion and preprocessing:} normalize raw events, deduplicate records, and assign region-time buckets.
\item \textbf{Signal extraction:} compute linguistic and temporal indicators from post streams.
\item \textbf{Credibility modeling:} estimate whether the observed regional signal is reliable.
\item \textbf{Threat filtering:} suppress non-organic amplification and media-driven artifacts.
\item \textbf{Decision engine:} map the final score to operational escalation states.
\item \textbf{Human review and audit:} retain accountability for any high-impact action.
\end{itemize}

Each stage produces intermediate values that can be inspected, explained, and debugged.

\subsection{Data Flow}

The system turns raw social media events into region-level alerts through a staged feature and scoring pipeline. The goal is to make sure every escalation can be traced back to observed evidence, a credibility adjustment, and a threat-suppression term.

\begin{figure}[htbp]
\centering
\includegraphics[width=0.9\textwidth]{dataflow.png}
\caption{Data flow from raw social media posts to region-level alerts.}
\label{fig:dataflow}
\end{figure}

\vspace{0.5cm}

As shown in Figure \ref{fig:dataflow}, the raw event stream is first cleaned into structured posts, then converted into feature vectors, then adjusted by the credibility layer, then filtered for adversarial behavior, and finally mapped into escalation decisions.

\subsection{Interfaces}

The ingestion stage exposes cleaned event records with region, timestamp, text, and provenance metadata. The signal extraction stage emits linguistic, behavioral, and temporal feature vectors. The credibility layer returns a bounded reliability score. The threat filter returns a bounded suppression score. The decision engine consumes these outputs together and emits a categorical action state with audit metadata.

\section{Detailed Design}

\subsection{Modules}

The proposal uses a modular design so that each subsystem can be tested and improved independently. The most important modules are:

\begin{itemize}
\item \textbf{Ingestion module:} validates and standardizes raw data.
\item \textbf{Signal module:} extracts crisis-related features from language and behavior.
\item \textbf{Credibility module:} penalizes unstable or weakly supported signals.
\item \textbf{Threat module:} detects lexical repetition, abnormal URL density, and coordinated posting.
\item \textbf{Decision module:} performs confidence-gated escalation.
\end{itemize}

Figure \ref{fig:modules} shows how the core computation combines credibility, base crisis probability, and threat suppression.

\begin{figure}[htbp]
\centering
\includegraphics[width=0.78\textwidth]{modules.png}
\caption{Module-level view of credible signal computation.}
\label{fig:modules}
\end{figure}

\vspace{0.5cm}

\subsection{Algorithms}

The final score is defined as:

\[
\text{credible\_signal} = \chi \times P(\text{crisis} \mid s) \times (1 - \text{threat\_score})
\]

where $P(\text{crisis} \mid s)$ is the base crisis likelihood given extracted features $s$, $\chi$ is the credibility multiplier, and $\text{threat\_score}$ is the adversarial suppression term. This is useful because it separates strong evidence from trustworthy evidence.

The credibility multiplier is computed as:

\[
\chi = \text{sample\_stability} \times \text{spatial\_support} \times \text{temporal\_consistency}
\]

\[
\text{sample\_stability} = \min\left(1, \sqrt{\frac{N}{k}}\right)
\]

This term penalizes low-volume regions, rewards agreement across neighboring regions, and reduces the impact of short-lived spikes that differ sharply from historical behavior.

\subsection{Decision Logic}

The decision engine maps the final credible-signal score into simple operational states such as log-only, monitor, human review, and priority alert. Thresholds are intentionally conservative, and high-priority actions depend on both score magnitude and supporting context.

\begin{figure}[htbp]
\centering
\includegraphics[width=0.62\textwidth]{decision_flow.png}
\caption{Escalation decision flow used by the confidence-gated decision engine.}
\label{fig:decision}
\end{figure}

\vspace{0.5cm}

As shown in Figure \ref{fig:decision}, escalation is not a single threshold check. The system first filters low-signal events, then separates monitoring cases from human-review cases, and only routes to priority alert when there is extra cluster evidence.

\section{Implementation Plan}

The implementation order is intentionally simple so that each stage can be tested before the next one is added.

\subsection*{Phase 1: Data and Signal Pipeline}
\begin{itemize}
\item Finalize replayable input schema and preprocessing logic.
\item Implement region mapping, deduplication, and event normalization.
\item Add sentiment, lexicon, volume, and velocity feature extraction.
\end{itemize}

\subsection*{Phase 2: Credibility and Threat Layers}
\begin{itemize}
\item Implement sample stability, spatial support, and temporal consistency.
\item Add lexical entropy, URL density, and coordinated-velocity detectors.
\item Integrate both layers into a single credible-signal computation path.
\end{itemize}

\subsection*{Phase 3: Decision, Evaluation, and Governance}
\begin{itemize}
\item Implement escalation thresholds and human-review routing.
\item Add audit logging and explanation payloads.
\item Run scenario-based evaluation and ablation studies.
\item Harden the pipeline for demonstration and documentation.
\end{itemize}

\section{Timeline}

The timeline below is structured as a 12-week GSoC execution plan.

\begin{table}[htbp]
\centering
\caption{Planned weekly execution schedule}
\label{tab:timeline}
\begin{tabularx}{\textwidth}{@{}c X@{}}
\toprule
\textbf{Week} & \textbf{Work} \\
\midrule
1--2 & Repository setup, replay loader, preprocessing scaffold, and region-mapping strategy. \\
3 & Signal extraction for sentiment, lexicon hits, activity volume, and rate-of-change features. \\
4 & Sample-stability and temporal-consistency implementation. \\
5 & Spatial support logic and neighborhood-level aggregation. \\
6 & Threat-filter indicators for entropy, URL density, and coordinated posting. \\
7 & Integrated credible-signal computation and thresholded decision engine. \\
8 & Human-review routing, logging, and API integration. \\
9 & Sparse-data and bot-amplification scenario generation. \\
10 & Media-spike and organic-crisis evaluation runs. \\
11 & Calibration checks, threshold tuning, and ablation studies. \\
12 & Final hardening, documentation, and submission packaging. \\
\bottomrule
\end{tabularx}
\end{table}

\section{Evaluation Plan}

Evaluation focuses on reliability rather than only model accuracy. The system must stay conservative during adversarial or informational spikes while still surfacing credible, organically emerging crisis patterns.

\begin{figure}[htbp]
\centering
\includegraphics[width=0.9\textwidth]{evaluation.png}
\caption{Evaluation scenarios under sparse, adversarial, media-driven, and organic conditions.}
\label{fig:evaluation}
\end{figure}

\vspace{0.5cm}

Figure \ref{fig:evaluation} illustrates the four primary evaluation scenarios:

\begin{itemize}
\item \textbf{Sparse data:} verify that low-volume regions are damped rather than over-escalated.
\item \textbf{Bot attack:} verify that coordinated amplification is suppressed.
\item \textbf{Media spike:} verify that informational bursts do not trigger crisis escalation.
\item \textbf{Organic crisis:} verify that multi-signal, persistent events are surfaced with high confidence.
\end{itemize}

Primary metrics are precision at a fixed recall point, false positive rate, calibration error, and temporal stability across replay runs.

\section{Challenges and Risks}

The main technical risks are false escalation, adversarial manipulation, sparse-data instability, and miscalibration. Each one has a direct mitigation in the design.

\begin{itemize}
\item \textbf{False escalation:} mitigated through credibility gating and mandatory human review.
\item \textbf{Adversarial manipulation:} mitigated through entropy, URL-density, and velocity filters.
\item \textbf{Sparse-data instability:} mitigated through sample-size penalties and minimum-support thresholds.
\item \textbf{Miscalibration:} mitigated through explicit calibration evaluation and threshold tuning.
\end{itemize}

\begin{figure}[htbp]
\centering
\includegraphics[width=0.88\textwidth]{risk_flow.png}
\caption{Risk-to-mitigation mapping for the proposed crisis-detection pipeline.}
\label{fig:risk-flow}
\end{figure}

\vspace{0.5cm}

Figure \ref{fig:risk-flow} summarizes the direct mapping between major failure modes and the system mechanisms intended to control them.

\section{Prior Work}

I have already built a prototype codebase around this idea, including backend services for ingestion, enrichment, scoring, alerts, and API routing, along with an initial frontend dashboard for monitoring outputs. This reduces startup risk and provides a concrete base for integrating the credibility layer, threat filtering, and governance hooks proposed here.

The current work also includes lexicon-based crisis indicators, replayable pipelines, test scaffolding, and early architecture diagrams. This means the GSoC work can focus on credible-signal modeling, evaluation, and system hardening rather than starting from scratch.

\section{Conclusion}

This proposal addresses a clear weakness in current AI-based mental-health monitoring systems: they may detect signal strength, but they often do not show whether that signal is reliable enough to act on. The proposed credible signal detection system adds explicit reliability modeling, adversarial suppression, and human-in-the-loop escalation to close that gap.

The result is a more deployable and auditable pipeline for public-health monitoring. Its value is not only in improving raw prediction quality, but in producing alerts that remain stable, interpretable, and trustworthy in real-world conditions.

\end{document}
